\documentclass[border=4pt]{standalone}
\usepackage{amsmath,amssymb,mathtools,xcolor,varwidth}
\definecolor{proj}{HTML}{B22222}
\definecolor{resist}{HTML}{4682B4}
\definecolor{gravity}{HTML}{6A0DAD}
\begin{document}
\begin{varwidth}{28em}
\large\bfseries
\textcolor{proj}{Projectiles} persevere in their motions, so far as they are not retarded by the \textcolor{resist}{resistance of the air}, \     or impelled downwards by the \textcolor{gravity}{force of gravity}. A stone or a cannonball, once flung, would fly straight \     on forever were these two forces absent — shot fast enough horizontally, it would circle the whole Earth. \     What we see as a falling arc is only the \textcolor{proj}{projectile's} straight perseverance bent by $\vec{a} = \vec{g} + \tfrac{\vec{F}_{\text{air}}}{m}$.
\end{varwidth}
\end{document}
